\documentclass[11pt,a4paper]{article}
\usepackage{fontspec}
\usepackage[margin=1in]{geometry}
\usepackage[unicode,hidelinks]{hyperref}
\setlength{\parindent}{0pt}
\setlength{\parskip}{.6em}
\emergencystretch=3em
\newcommand{\volumeentry}[3]{\par\begingroup\parskip=.7em
  \noindent\href{Consciousness_Trilogy_Complete_Collection.pdf\#page=#3}{\textbf{#1}\quad #2\nobreak\dotfill\nobreak\mbox{#3}}\par\endgroup}
\hypersetup{pdftitle={Consciousness Trilogy},pdfauthor={Sidong Liu}}
% COVER SUBTITLE: Geometry, perspective and reflective organization
\begin{document}
\begin{titlepage}
\centering
\vspace*{1.6cm}
{\Huge\bfseries Consciousness Trilogy\par}
\vspace{1cm}
{\Large Geometry, perspective and reflective organization\par}
\vfill
{\Large Complete Collected Volume\par}
\vspace{1.6cm}
{\Large Sidong Liu, PhD\par}
\vspace{.3cm}
{iBioStratix Ltd\par}
\href{mailto:sidongliu@hotmail.com}{\texttt{sidongliu@hotmail.com}}\par
\vspace{1cm}
{\large October 2026\par}
\end{titlepage}
\setcounter{page}{2}
\section*{Contents}
The numbers below refer to pages in this PDF. Each paper retains its own printed pagination, proofs, appendices and references.
\volumeentry{}{Introduction and reading map}{3}
\volumeentry{C-I}{Geometry Relative to a Cut: Selection, Generation and Modular Inseparability}{4}
\volumeentry{C-II}{Higher-Dimensional Consistency of Geometry and Modular Flow: Physical Source Identification and a Second-Order Extension}{28}
\volumeentry{C-III}{Finite Self-Models and the Explanatory Gap: What a Continuing Observer Can Know About Itself}{64}
\clearpage
\section*{Introduction}\label{collection:intro}
How does a world become available from a perspective, and what would connect that availability to first-person organization? These papers retain both the structural question and the explanatory demand. Geometry cannot reach a subject merely by being reconstructed by an external theorist.

I studies the correspondence between an internal standard Cut and modular geometry in its stated net class. It distinguishes a real loss of accessible distinctions from reversible recoding. II identifies physical source data and tests higher-dimensional geometric consistency; its common-source and preparation conditions keep the geometric test attached to matter rather than an arbitrary fit.

III joins finite self-description to historically formed access and causal self-diagnosis. One two-layer bridge asks first for an actual path from geometry through instruments and decoded records to usable reports and actions, then for independent evidence connecting reflective organization to experience. Competing mechanisms, sham interventions and predeclared markers are part of the proposal. Neither a Cut nor a successful calibration experiment alone establishes phenomenal presence. The functional tests distinguish display-only explanations, but not every general metacognitive controller; independent reflective evidence remains necessary.

\section*{Reading map}
\begin{description}
\item[C-I] Internal standard Cuts determine modular geometry in the stated class; the perspective argument distinguishes real access restrictions from recoding.
\item[C-II] Prepared source families and modular data test higher-dimensional geometry. Positive-energy Radon identifiability and common-source conditions keep the field-law claim conditional.
\item[C-III] Finite self-models have permission-dependent limits; one two-layer bridge connects internal geometric access to reflective organization and its independent empirical tests.
\end{description}
\end{document}
